\section{Results and Discussion}

\subsection{Headings}
Headings should be capitalized (\ie, nouns, verbs, and all other words except articles, prepositions, and conjunctions should be set with an initial capital) and should, with the exception of the title, be aligned to the left.
Only the first two levels of section headings should be numbered, as shown in \cref{tab:headings}.
The respective font sizes are also given in \cref{tab:headings}. 
Kindly refrain from using ``0'' when numbering your section headings.
Words joined by a hyphen are subject to a special rule. 
If the first word can stand alone, the second word should be capitalized.

\begin{table}[tb]
  \caption{Font sizes of headings. 
    Table captions should always be positioned \emph{above} the tables.
  }
  \label{tab:headings}
  \centering
  \begin{tabular}{@{}lll@{}}
    \toprule
    Heading level & Example & Font size and style\\
    \midrule
    Title (centered)  & {\Large\bf Lecture Notes \dots} & 14 point, bold\\
    1st-level heading & {\large\bf 1 Introduction} & 12 point, bold\\
    2nd-level heading & {\bf 2.1 Printing Area} & 10 point, bold\\
    3rd-level heading & {\bf Headings.} Text follows \dots & 10 point, bold\\
    4th-level heading & {\it Remark.} Text follows \dots & 10 point, italic\\
  \bottomrule
  \end{tabular}
\end{table}

Here are some examples of headings: 
``Criteria to Disprove Context-Freeness of Collage Languages'', ``On Correcting the Intrusion of Tracing Non-deterministic Programs by Software'', ``A User-Friendly and Extendable Data Distribution System'', ``Multi-flip Networks: Parallelizing GenSAT'', ``Self-determinations of Man''.


\subsection{Figures}
\label{sect:figures}
For \LaTeX{} users, we recommend integrating figures in your paper using the package \texttt{graphicx}.

It is essential that all illustrations are clear and legible. 
Vector graphics (rather than rasterized images) should be used for diagrams and schemas whenever possible. 
Please check that the lines in line drawings are not interrupted and have a constant width. 
Line drawings are to have a resolution of at least 800 dpi (preferably 1200 dpi).
Grids and details within figures must be clearly legible and may not be written one on top of the other. 
The lettering in figures should not use font sizes smaller than 6\:pt ($\sim$2\:mm character height). 

Figures should be numbered and should have a caption, which should always be positioned \emph{under} the figures, in contrast to the caption belonging to a table, which should always appear \emph{above} the table.
Figures and Tables should be cross-referred in the text.

If they are short, they are centered between the margins (\cf \cref{fig:short}). 
Longer captions, covering more than one line, are justified (\cref{fig:example} shows an example). 
Captions that do not constitute a full sentence, do not have a period.


\begin{figure}[tb]
  \centering
  \includegraphics[height=6.5cm]{eijkel2}
  \caption{One kernel at $x_s$ (\emph{dotted kernel}) or two kernels at $x_i$ and $x_j$ (\emph{left and right}) lead to the same summed estimate at $x_s$.
    This shows a figure consisting of different types of lines.
    Elements of the figure described in the caption should be set in italics, in parentheses, as shown in this sample caption. 
    The last sentence of a figure caption should generally end with a full stop, except when the caption is not a full sentence.
  }
  \label{fig:example}
\end{figure}

\begin{figure}[tb]
  \centering
  \begin{subfigure}{0.68\linewidth}
    \fbox{\rule{0pt}{0.5in} \rule{.9\linewidth}{0pt}}
    \caption{An example of a subfigure}
    \label{fig:short-a}
  \end{subfigure}
  \hfill
  \begin{subfigure}{0.28\linewidth}
    \fbox{\rule{0pt}{0.5in} \rule{.9\linewidth}{0pt}}
    \caption{Another example of a subfigure}
    \label{fig:short-b}
  \end{subfigure}
  \caption{Centered, short example caption}
  \label{fig:short}
\end{figure}

If possible (\eg, if you use \LaTeX) please define figures as floating objects. 
\LaTeX{} users, please avoid using the location parameter ``h'' for ``here''. 
If you have to insert a pagebreak before a figure, please ensure that the previous page is completely filled.


\subsection{Formulas}
Displayed equations or formulas are centered and set on a separate line (with an extra line or half line space above and below). 
Equations should be numbered for reference. 
The numbers should be consecutive within the contribution, with numbers enclosed in parentheses and set on the right margin.
For example,
\begin{align}
  \psi (u) & = \int_{0}^{T} \left[\frac{1}{2}
  \left(\Lambda_{0}^{-1} u,u\right) + N^{\ast} (-u)\right] \text{d}t \; \\
& = 0
\end{align}
and 
\begin{equation}
  E = m\cdot c^2.
  \label{eq:important}
\end{equation}
Please do not include section counters in the numbering.

Numbering equations makes reviewing more efficient, because reviewers can refer to a line on a page.  
It is important for readers to be able to refer to any particular equation.
Just because you did not refer to it in the text does not mean some future reader might not need to refer to it.
It is cumbersome to have to use circumlocutions like ``the equation second from the top of page 3''.
(Note that the ruler will not be present in the final copy, so is not an alternative to equation numbers).
All authors will benefit from reading Mermin's description of how to write mathematics:
\url{https://doi.org/10.1063/1.2811173}.

% No color equations in Springer publications.
Equations should never be in color and should be punctuated in the same way as ordinary text.
They should not be pasted in as figures.


\subsubsection{Lemmas, Propositions, and Theorems.}
The numbers accorded to lemmas, propositions, and theorems, \etc should appear in consecutive order, starting with Lemma 1. 
Please do not include section counters in the numbering like ``Theorem 1.1''.


\subsection{Footnotes.}
The superscript numeral used to refer to a footnote appears in the text either directly after the word to be discussed or -- in relation to a phrase or a sentence -- following the punctuation mark (comma, semicolon, or period).%
\footnote{The footnote numeral is set flush left and the text follows with the usual word spacing. 
  Second and subsequent lines are indented. 
}
For remarks pertaining to the title or the authors' names, in the header of a paper, symbols should be used instead of a number.
Please note that no footnotes may be included in the abstract.


\subsection{Cross References}
For the benefit of author(s) and readers, please use the
\begin{verbatim}
  \cref{...}
\end{verbatim}
command for cross-referencing to figures, tables, equations, or sections.
This will automatically insert the appropriate label alongside the cross reference as in this example:
\begin{quotation}
  To see how our method outperforms previous work, please see \cref{fig:example} and \cref{tab:headings}.
  It is also possible to refer to multiple targets as once, \eg~to \cref{fig:example,fig:short-a}.
  You may also return to \cref{sec:intro} or look at \cref{eq:important}.
\end{quotation}
If you do not wish to abbreviate the label, for example at the beginning of the sentence, you can use the
\begin{verbatim}
  \Cref{...}
\end{verbatim}
command. Here is an example:
\begin{quotation}
  \Cref{fig:example} is also quite important.
\end{quotation}


\subsection{Program Code}
Program listings or program commands in the text are normally set in typewriter font (\eg, \texttt{printf("Hello world!\textbackslash{}n");}).


\subsection{Citations}
Arabic numbers are used for citation, which is sequential either by order of citation or by alphabetical order of the references, depending on which sequence is used in the list of references. 
The reference numbers are given in brackets and are not superscript.
Please observe the following guidelines:
\begin{itemize}
\item Single citation: \cite{Anonymous24}
\item Multiple citation: \cite{Alpher02,Alpher03,Alpher05,Anonymous24b,Anonymous24}. 
  The numbers should be listed in numerical order.
  If you use the template as advised, this will be taken care of automatically.
\item If an author's name is used in the text: Alpher \cite{Alpher02} was the first \ldots
\end{itemize}
Please write all references using the Latin alphabet. If the title of the book you are referring to is, \eg, in Russian or Chinese, then please write (in Russian) or (in Chinese) at the end of the transcript or translation of the title.
All references cited in the text should be in the list of references and vice versa.

References should be formatted with the official LNCS reference style.
The \LaTeX{} template already takes care of that through the use of the \texttt{splncs04.bst} Bib\TeX{} style file.
Springer strongly encourages you to include DOIs (Digital Object Identifiers) in your references (\cf \cite{ECCV2022}). 
The DOI is a unique code allotted by the publisher to each online paper or journal article. 
It provides a stable way of finding published papers and their metadata. 
The insertion of DOIs increases the overall length of the references section, but this should not concern you as the reference section is not counted toward the page limit.


\subsection{Miscellaneous}
Compare the following:
\begin{center}
  \begin{tabular}{ll}
    \verb'$conf_a$'          & $\qquad conf_a$ \\
    \verb'$\mathit{conf}_a$' & $\qquad \mathit{conf}_a$
  \end{tabular}
\end{center}
See The \TeX book, p.\ 165.

The space after \eg, meaning ``for example'', should not be a sentence-ending space.
So \eg is correct, \emph{e.g.} is not.
The provided \verb'\eg' macro takes care of this.

When citing a multi-author paper, you may save space by using ``et alia'', 
shortened to ``\etal'' (not ``{\em et.\ al.}'' as ``{\em et\hskip 0.1em}'' is a complete word).
If you use the \verb'\etal' macro provided, then you need not worry about double periods when used at the end of a sentence as in Alpher \etal.
However, use it only when there are three or more authors.
Thus, the following is correct:
   ``Frobnication has been trendy lately.
   It was introduced by Alpher~\cite{Alpher02}, and subsequently developed by
   Alpher and Fotheringham-Smythe~\cite{Alpher03}, and Alpher \etal~\cite{Alpher04}.''

This is incorrect: ``... subsequently developed by Alpher \etal~\cite{Alpher03} ...'' because reference~\cite{Alpher03} has just two authors.

